\subsection{Component detection and the meaning of ``useful''}

\begin{claim}{G.1}{Proved}{Locally constant coefficients detect components for every nontrivial derivator}
Let $\D$ be a derivator on all $\U$-small categories. If $\D(e)$ is
not equivalent to the terminal category, then
\[
 \Wlc\subseteq\W_0.
 \tag{G.1.1}
\]
If $\D(e)\simeq e$, then $\D(A)\simeq e$ for every $A$, and all
five classes consist of every functor. Thus the hypothesis is exact.
\end{claim}
\begin{proof}
There are two nonisomorphic objects $M,N\in\D(e)$: a category with
an initial object and only one isomorphism class of objects is
equivalent to $e$. The initial object exists by Der~1 and Der~3.
Suppose first that $Bf$ misses a component of $BB$. Use Der~1 to
choose two coefficients on $B$ constant with value $M$ on all other
components and respectively $M,N$ on the missed component. They are
locally constant, and their pullbacks are isomorphic. An equivalence
$f^*$ would make them isomorphic on $B$, contradicting evaluation
at a point of the missed component.

If two distinct components of $A$ map into the same component of
$B$, choose a locally constant coefficient on $A$ taking the values
$M,N$ on those components, and $M$ elsewhere. It cannot be a
restriction: any locally constant object on the connected component
of $B$ has isomorphic values at all its objects, along a finite zigzag
of arrows. This contradicts essential surjectivity. These arguments
also exclude $\varnothing\to B$ for nonempty $B$ and handle
arbitrarily many components without a finiteness assumption.

If $\D(e)\simeq e$, all arrows of $\D(A)$ are isomorphisms by
Der~2. The category $\D(A)$ has an initial object, from the left Kan
extension along $\varnothing\to A$; a groupoid with an initial
object is equivalent to $e$. Every restriction is then an
equivalence and every comparison is invertible.
\end{proof}

\begin{claim}{G.2}{Proved}{The exact obstruction for constant coefficients}
For a derivator, $\Wc\subseteq\W_0$ holds if and only if
$\D(e)$ is not a preorder up to equivalence. More explicitly, if
there exist two distinct parallel morphisms $\alpha,\beta:M\to N$
in $\D(e)$, then every constant-coefficient equivalence is bijective
on components. If $\D(e)$ is a nontrivial preorder up to equivalence,
then
\[
 \Wc=\{f:A\to B:\ A\text{ and }B\text{ are both empty or both nonempty}\}.
 \tag{G.2.1}
\]
In particular, Chapter X's component-detection expectation needs a
coefficient-separation hypothesis for $\Wc$, but only nontriviality
for $\Wlc$.
\end{claim}
\begin{proof}
Suppose $\alpha\ne\beta$. If a target component is missed, define
two morphisms $q^*M\to q^*N$ using $\alpha$ everywhere and changing
it to $\beta$ just on that component. Der~1 supplies these maps;
evaluation distinguishes them, while restriction identifies them.
If two source components coalesce, define a morphism $p^*M\to p^*N$
using different choices on them. Every morphism between constants
over $B$ induces, under the underlying diagram functor, a natural
transformation between constant presheaves. Its component morphisms
are equal along every arrow of $B$, hence equal throughout each
connected component. The chosen source map is therefore not in the
image of (A.1.1). Full bijectivity excludes both failures.

Now assume $\D(e)$ is thin. For nonempty $A$, choose $a:e\to A$.
The adjunction unit $N\to p_*p^*N$ and evaluation at $a$ give maps
in both directions between $N$ and $p_*p^*N$. In a thin category
these are mutually inverse. Thus
\[
 \Hom_{\D(A)}(p^*M,p^*N)
 \cong\Hom_{\D(e)}(M,p_*p^*N)
 \cong\Hom_{\D(e)}(M,N).
\]
Under these identifications every restriction between nonempty
shapes is the identity. On the empty shape every Hom set is a
singleton. In a nontrivial complete preorder the morphism from
terminal to initial is absent; that pair distinguishes
$\varnothing\to B$ for nonempty $B$. This proves (G.2.1).
The represented complete lattice $0<1$ is an explicit example:
constants see only emptiness, while its locally constant coefficient
category is $(0<1)^{\pi_0A}$ and detects $\W_0$ under equivalence.
\end{proof}

\subsection{Functoriality in the coefficient theory}

\begin{claim}{G.3}{Proved}{What an arbitrary morphism of derivators supplies}
A pseudonatural morphism $\Phi:\D\to\mathbf E$ restricts to
functors $\Phi_A:\LC_\D(A)\to\LC_{\mathbf E}(A)$, coherently
commuting with restriction. This alone imposes no inclusion in
either direction between $\W_\D^{\mathrm{lc}}$ and
$\W_{\mathbf E}^{\mathrm{lc}}$. Equivalence of derivators gives
equality of the five classes. Products give intersections of the
five classes.
\end{claim}
\begin{proof}
Pseudonaturality at objects and arrows identifies the underlying
presheaf of $\Phi_A F$ with the image under $\Phi_e$ of that of
$F$. Functors preserve isomorphisms, proving the first assertion.
There are morphisms in both directions between any derivator and
the terminal derivator: the unique functor in one direction and
the terminal-object section in the other. Restrictions preserve
terminal objects because they have left adjoints. The terminal
derivator has every functor in its class, while the derivator of
sets has $\W^{\mathrm{lc}}=\W_1$. Thus no monotonicity direction
is forced by the existence of an arbitrary morphism. An equivalence
of derivators identifies each of the specified comparisons.
For a product derivator, restrictions and all adjunction units and
counits are products; their invertibility and full faithfulness are
componentwise. Essential surjectivity is componentwise as well.
\end{proof}

\begin{claim}{G.4}{Proved}{Retracts and fully faithful comparisons}
If $\mathbf E$ is a pseudoretract of $\D$, namely if there are
morphisms $\Phi:\D\to\mathbf E$ and $\Psi:\mathbf E\to\D$
with $\Phi\Psi\cong\id$, then
\[
 \W_\D^{\mathrm{lc}}\subseteq\W_{\mathbf E}^{\mathrm{lc}},
 \qquad \W_\D^1\subseteq\W_{\mathbf E}^1,
 \qquad \W_\D\subseteq\W_{\mathbf E}.
 \tag{G.4.1}
\]
If instead every $\Phi_A$ is fully faithful, then
$\W_{\mathbf E}^1\subseteq\W_\D^1$ and
$\W_{\mathbf E}\subseteq\W_\D$. To infer the analogous assertion
for $2$-completeness one additionally needs descent of the essential
image: whenever $F\in\LC_\D(A)$ and
$G\in\LC_{\mathbf E}(B)$ satisfy
$f^*G\cong\Phi_A F$, an object $G$ with this property must be
choosable in the essential image of $\Phi_B$.
\end{claim}
\begin{proof}
The restriction functor for $\mathbf E$ is a retract, up to the
specified coherent isomorphism, of that for $\D$. Equivalences
and fully faithful functors are closed under these retracts by the
Hom-set and essential-surjectivity argument of A.4. Constants are
preserved by both morphisms, and (A.1.1) for $\mathbf E$ is a
retract of the corresponding map for $\D$ with coefficients
$\Psi_eM,\Psi_eN$. This proves (G.4.1).
For a fully faithful $\Phi$, the two vertical maps in every Hom
comparison square are bijections, so the asserted implications
follow. Essential surjectivity has a separate quantifier over
source coefficients; the last displayed condition supplies exactly
the missing lift of the descended coefficient. A preimage under
$\Phi_B$ is locally constant because the fully faithful $\Phi_e$
reflects isomorphisms on all underlying transition arrows. No conservativity
claim is substituted for that existence condition.
\end{proof}

\begin{claim}{G.5}{Proved}{Change of scalars and the Chapter X assignment}
Let $k\to l$ be a homomorphism of associative unital rings. Derived
extension of scalars gives
\[
 \W_{D(k)}^{\mathrm{lc}}\subseteq\W_{D(l)}^{\mathrm{lc}},
 \qquad \W_{D(k)}\subseteq\W_{D(l)}.
 \tag{G.5.1}
\]
The first inclusion is equality if derived extension of scalars
reflects equivalences and the canonical loop-algebra comparison
can be reflected by it; in particular this holds for a faithfully
flat map of nonzero commutative rings. Thus $\D\mapsto\Wlc$ is
well defined on equivalence classes of derivators and has precise
functoriality under additional hypotheses, but is not monotone on
all morphisms of derivators.
\end{claim}
\begin{proof}
The comparison of loop chains for $l$ is obtained from the
comparison for $k$ by derived scalar extension; the same statement
holds for ordinary singular chains. C.7 therefore proves (G.5.1)
when the rings are nonzero. A zero target ring gives the trivial
coefficient theory, and a unital map from the zero ring has zero
target, so the exceptional cases satisfy the same assertion.
For faithfully flat commutative extension, a complex is acyclic
if and only if its scalar extension is acyclic. The loop criterion
then reflects membership, including its component and
fundamental-group conditions. The final assertions follow from
G.3 and G.4.
\end{proof}

\begin{claim}{G.6}{Open}{A functorial higher-coefficient construction on bare derivators}
For an arbitrary derivator $\D$ with no specified monoidal
enhancement, this work does not construct a functorially assigned
derivator $\D'$ with $\W_{\D'}=\Wlc$, nor an iteratable
$2$-derivator construction of that kind. The requested general
statement is retained with those quantifiers. A chosen presentable
linear enhancement has the positive realization in E.2; a specified
accessible space localization has the conditional realization in E.7.
The missing data in the general proposal are a category of
$\D$-linear categories, its morphisms and coherent $2$-morphisms,
and Kan extensions on the enriched shapes required by Di Vittorio.
The full relative Theorem A for arbitrary derivators is proved
in B.3. That localizer theorem by itself does not specify an
iteratable coefficient category or the required enriched shapes.
\end{claim}
